% ID: SEL-D0047-OB-E03
% ESTADO: cotejada. Segunda pasada visual y borrador compilado revisados.
% Fuente: Selectividad/sel_2014_matematicas_aplicadas/Matematicas_CCSS_SEPT.pdf
% SHA-256: ee5bd8e23465d55ed58339e7d9adf3d58423ac2966b53d92e39a2f9b35558ab1
% Página: 2; opción B; ejercicio original 3.

En una tienda de complementos disponen de 100 bolsos, de los cuales 80 son de una conocida marca y 20 son imitaciones casi perfectas de dicha marca. Una inspección encarga a un experto el peritaje de los bolsos de la tienda. Se sabe que este experto acierta en el 95\% de sus peritajes cuando el bolso es auténtico y que detecta el 98\% de las imitaciones. Se elige, al azar, un bolso para su examen:
\begin{enumerate}[label=\alph*)]
\item \textbf{(1.25 puntos)} Calcule la probabilidad de que el experto acierte en su dictamen sobre ese bolso.
\item \textbf{(1.25 puntos)} Si el experto no ha acertado en su peritaje, calcule la probabilidad de que el bolso sea auténtico.
\end{enumerate}
